\section{Engineering Interpretation}
\label{sec:limits}
The results support an integrated systems contribution: industrial agent
execution can preserve authority, production context, command confirmation,
and process observations as distinct, inspectable evidence. Admission
ablation explains the contribution of the pre-write interval, the
publish-only challenge tests interpretation of command receipts, and the
paired trajectories connect interventions to observed targets. Their
combination makes an agent's completed task interpretable in terms of the
industrial operations it performed.

This evidence structure also explains the division between strategy and
platform. An LLM, scripted policy, or deterministic controller supplies a
candidate through its authorized interface. The platform applies the
execution contract and returns evidence with stable semantics. Likewise,
process reuse changes the signal map, objective, and guard definitions
while retaining production context and execution records. Q6 closes this
chain at the operational level: the retained recipe version follows
observed acceptance of an approved trial sequence.

Human participation is represented through explicit authority and
attributed decisions. This supports inspection of what an approver
authorized and what evidence the proposer later used. Automation research
distinguishes such workflow structure from operator comprehension,
complacency, and allocation of responsibility
\cite{parasuramanmanzey2010complacency,parasuraman2000model,bainbridge1983ironies}.
The server-side boundary resolves execution authority independently of
prompts and retrieved documents \cite{greshake2023injection,liu2024pinjection}.
Deployment combines this boundary with credential management, endpoint
security, and IT/OT separation \cite{iec62443-3-3}. The supervisory layer
coexists with independently engineered PLC interlocks and safety
instrumented systems; IEC~61511 governs the latter's process-safety
lifecycle \cite{iec61511}.

\noindent\textit{Experimental scope.}
The process evidence uses accelerated simulators, scripted approvals, and
sequential trajectories. Q1--Q4 cover the historical direct-node interface;
current profile and composition tests cover software behavior; Q6 covers
one recipe trajectory. A complete current AML-exploration control--response
trial is not archived. Q6's observed guards, snapshots, and cadence settings
are specified in Section~\ref{sec:recipeexperiment}. These evidence units
establish the exercised
mechanisms without estimating comparative optimization benefit, fault
frequency, human performance, hard real-time behavior, or physical transfer.
Matched initial states, repeated independent runs, and a fixed-policy or
SCADA comparator would support those further estimates. Established process
benchmarks such as Tennessee Eastman \cite{downs1993tep} provide an
additional setting for that evaluation.

\section{Conclusion}
\label{sec:conclusion}
The framework connects industrial agent teams to production-scoped,
governed execution while keeping proposal strategy replaceable. Semantic
nodes, capability/profile enforcement, source-specific admission, approval,
graded confirmation, and distinct node/recipe records form an inspectable
chain from a task to commanded controls, observed outputs, and retained
production parameters. Protocol and admission experiments, controller and
process missions, repeated regression execution, and the paired
injection-molding trajectory exercise complementary parts of that chain.
The contribution is a common execution and evidence contract through which
industrial agent reasoning can be connected to accountable supervisory
operation.

\noindent\textit{Data availability and AI assistance.}
Run archives, original Q2 engine logs, replay entries, configuration notes,
and claim-to-source mappings accompany the separate anonymous experimental
evidence bundle. The manuscript source package contains plotting sources
and figure mappings. An OpenAI language model assisted
with manuscript organization, language revision, and plotting-code
preparation. The authors remain responsible for the final submission.
